\section{Uniform amplification of commutators}
\label{sec:amplification}

The finite-factor argument will first produce an isomorphism that is close
to the given representation in operator norm. To implement it by a small
unitary, we need the same control at every matrix level. The following
estimate supplies that passage.

\begin{theorem}[Uniform amplification]\label{thm:amplification}
There is an absolute constant \(C_{\mathrm{amp}}\) such that, for every
type \(\II_1\) factor \(M\subseteq\B(H)\) on a separable Hilbert space and
every selfadjoint \(h\in\B(H)\),
\[
 \norm{\ad_h|_M}_{\cb}\le
 C_{\mathrm{amp}}\norm{\ad_h|_M}.
\]
The constant is independent of \(M\), its representation, and \(H\).
\end{theorem}

We first control columns by the cyclic similarity theorem. A free-word
estimate then permits us to apply this column bound to randomized matrix
coefficients. The randomization also recovers a fixed positive multiple
of the identity in a weak limit; this will detect every matrix
commutator. The auxiliary matrix sizes depend on the matrix level being
tested, but every estimate that survives to the conclusion is absolute.

\subsection{Columns and trace limits}

For \(a_1,\ldots,a_l\in\B(H)\), write
\(\norm{(a_i)_i}_{\mathrm{col}}=\norm{\sum_i a_i^*a_i}^{1/2}\).

\begin{lemma}[Column estimate]\label{lem:column-commutator}
Let \(A\subseteq\B(H)\) be a unital \(C^*\)-algebra, let \(h\in\B(H)\),
and put \(\delta=\norm{\ad_h|_A}\). Then
\[
 \norm{([h,a_i])_i}_{\mathrm{col}}
 \le 8\delta\,\norm{(a_i)_i}_{\mathrm{col}}
 \qquad(a_1,\ldots,a_l\in A).
\]
\end{lemma}

\begin{proof}
There is nothing to prove if \(\delta=0\). Otherwise
we use the triangular homomorphism associated with a derivation,
as in \cite[p.~216]{Haagerup1983}:
\[
 \Phi(a)=
 \begin{pmatrix}a&\delta^{-1}[h,a]\\0&a\end{pmatrix}
 \quad(a\in A).
\]
This is a unital homomorphism with \(\norm{\Phi}\le2\). For \(\xi\in H\),
the closed subspace
\[
 K_\xi=\overline{\{\Phi(a)(0,\xi):a\in A\}}
 \subseteq H\oplus H
\]
is invariant for \(\Phi(A)\), and the restricted homomorphism is cyclic.
Haagerup's cyclic similarity theorem gives
\(\norm{\Phi|_{K_\xi}}_{\cb}\le\norm{\Phi|_{K_\xi}}^3\le8\);
see \cite[Theorem~1.1 and its proof, pp.~217, 225]{Haagerup1983}.
Apply its column amplification to \((0,\xi)\in K_\xi\), and retain the
upper component of each output. This gives
\[
 \left(\sum_i\norm{[h,a_i]\xi}^2\right)^{1/2}
 \le8\delta\,\norm{(a_i)_i}_{\mathrm{col}}\norm{\xi}.
\]
Taking the supremum over unit vectors proves the claim.
\end{proof}

We use normalized traces on finite factors and the standard comparison
and subdivision of projections by trace
\cite[III.1.7.9--10, p.~257; III.2.5.4(iii) and Theorem~III.2.5.7,
p.~279]{BlackadarOperatorAlgebras}.
For a finite von Neumann
algebra \((P,\tau)\), put
\(\norm{x}_2=\tau(x^*x)^{1/2}\) and
\(\norm{x}_1=\tau(|x|)\). Fix a free ultrafilter \(\omega\) on
\(\N\). The tracial ultrapower \(P^\omega\) is the quotient of the
bounded sequences in \(P\) by those with
\(\lim_{j\to\omega}\norm{x_j}_2=0\); its trace is
\(\tau_\omega((x_j)_\omega)=\lim_{j\to\omega}\tau(x_j)\).
We identify \(P\) with its constant sequences.

\begin{lemma}[Transfer from the trace to a normal representation]
\label{lem:trace-transfer}
Let \(P\subseteq\B(K)\) be a finite von Neumann algebra with faithful
normal trace \(\tau\).
\begin{enumerate}[label=\textup{(\roman*)}]
\item If \((x_j)\) is bounded and
\(\lim_{j\to\omega}\norm{x_j}_2=0\), then
\(\lim_{j\to\omega}\norm{x_j\xi}=0\) for every \(\xi\in K\).
\item If \(x=(x_j)_\omega\in P^\omega\), then the weak-operator
ultralimit of \(x_j\) in this representation is \(E_P(x)\), where
\(E_P:P^\omega\to P\) is the trace-preserving conditional expectation.
\end{enumerate}
\end{lemma}

\begin{proof}
The functionals \(a\mapsto\tau(ba)\), \(b\in P\), are norm dense in
the predual \(P_*\). Indeed, their annihilator in
\((P_*)^*=P\) is zero, by faithfulness of the trace, so
Hahn--Banach gives the density.
For a normal vector functional \(\varphi\), approximate \(\varphi\)
in norm by one of these functionals. If \(\norm{x_j}\le K_0\), then
\[
 |\varphi(x_j^*x_j)|
 \le K_0^2\norm{\varphi-\tau(b\,\cdot)}
       +\norm{b}\norm{x_j}_2^2.
\]
This proves (i). For (ii), the defining identity of the conditional
expectation gives
\[
 \lim_{j\to\omega}\tau(bx_j)
 =\tau_\omega(bx)=\tau(bE_P(x))\qquad(b\in P).
\]
The same predual-density argument, using boundedness, extends the
identity to every normal functional.
\end{proof}

\subsection{Free projection words}

The following estimate controls sums of projection words whose labels have
the same multiplicities. In the later randomization, these sums will give
coefficients to which we apply the column estimate separately. When the
squared estimates are added, each diagonal vector component is counted
once for every multiset containing its label. We therefore need squared
word-norm bounds with a uniformly bounded sum through each fixed label.

\begin{lemma}[Free projection words]\label{lem:free-words}
For every integer $r\geq 1$ there are constants $C_r>0$ and an integer
$d_r\geq r$ with the following property.  Let $d\geq d_r$, and let
$(\mathcal A,\tau)$ be a von Neumann algebra with a faithful normal
tracial state.  Suppose that $A_1,\ldots,A_r\subset\mathcal A$ are
unital copies of $M_d(\C)$ which are scalar free: the trace of a product
of centered letters from successively different copies is zero.
For each $j$, let $(e_a^{(j)})_{a=1}^d$ be the diagonal minimal
projections of $A_j$.

For a multiset $W$ of cardinality $r$ drawn from $\{1,\ldots,d\}$, put
\[
 B_W=\sum_{\substack{(a_1,\ldots,a_r)\in\{1,\ldots,d\}^r\\
                     \text{with multiset }W}}
       e_{a_r}^{(r)}\cdots e_{a_1}^{(1)}.
\]
Then $\norm{B_W}\leq\beta_W$, where
\begin{equation}\label{eq:fw-beta}
 \beta_W=
 \begin{cases}
  7\sqrt{(r-1)!}\,d^{-(r-1)/2},
       &\text{if all labels in $W$ are distinct},\\
  C_r d^{-(r-1)/2},&\text{otherwise}.
 \end{cases}
\end{equation}
The constants and threshold depend only on $r$, not on the ambient
algebra, the free copies, or $W$.
\end{lemma}

Every word in $B_W$ ends with a projection $e_a^{(1)}$ whose label
occurs in $W$. Thus $B_W$ has right support in the sum of those
projections from the first copy, at most $r$ of them. This support
bound will let us reduce each coefficient to a column, once the
auxiliary matrix dimension is large enough.

\begin{proof}
We estimate the centered words first, then restore their scalar parts.
The decomposition below uses the creation, annihilation, and diagonal
actions in the free-product Hilbert space, as in the framework developed by
Voiculescu \cite{Voiculescu1985} and the reduced-free-product inequalities of
Ricard and Xu \cite[Section~1 and Fact~2.6]{RicardXu2006}.
The corresponding bound for words of fixed length in free groups is
Haagerup's inequality \cite[Lemma~1.4]{Haagerup1979}.
Here the prescribed label multiplicities require the additional tuple-matrix
and square-summation estimates proved below.

\paragraph{Reduced words and their cuts.}
We first describe the Hilbert space on which the word norm will be
estimated.  Replace $\mathcal A$ by $W^*(A_1,\ldots,A_r)$ and use its
faithful left representation on $L^2(\mathcal A,\tau)$.  Set
\[
 H_j^\circ=L^2(A_j,\tau)\ominus\C1.
\]
Freeness identifies this representation space with
\begin{equation}\label{eq:fw-fock}
 \mathcal F=\C\Omega\ \oplus\!
 \bigoplus_{\substack{\ell\geq1,\ i_1,\ldots,i_\ell\in\{1,\ldots,r\}\\
                      i_t\ne i_{t+1}}}
 H_{i_1}^\circ\otimes\cdots\otimes H_{i_\ell}^\circ.
\end{equation}
Indeed, map a tensor to the corresponding product applied to $1$.
Its inner products are the tensor inner products: in the trace of the
adjoint of one reduced word times another, pair the adjacent letters
at the junction, split their product into its trace and centered part,
and iterate.  Freeness eliminates the centered terms until the copy
sequences agree and all letters have been paired.  Distinct summands
are orthogonal.  Splitting each letter into its scalar and centered
parts shows that reduced words span the algebra generated by the
copies, hence a dense subspace of $L^2(\mathcal A,\tau)$.

Write $h_a^{(j)}=e_a^{(j)}-d^{-1}1$.  Define the label map
\[
 D_j:\ell^2(\{1,\ldots,d\})\longrightarrow H_j^\circ,
 \qquad D_j\delta_a=h_a^{(j)}.
\]
The normalized matrix trace gives
\begin{equation}\label{eq:fw-label-map}
 D_j^*D_j=\frac1d I-\frac1{d^2}\mathbf1\mathbf1^*,
 \qquad \norm{D_j}\leq d^{-1/2},
\end{equation}
where $\mathbf1$ is the vector all of whose entries are one.

Left multiplication by $h_a^{(j)}$ creates the first letter
$h_a^{(j)}$ on a reduced tensor not beginning in copy $j$.
On a tensor $b\otimes\xi$ beginning in that copy, it has two parts:
annihilation, equal to $\tau(h_a^{(j)}b)\xi$, and preservation of the
first leg, with operator
\begin{equation}\label{eq:fw-preservation-operator}
 K_a^{(j)}
 =P_j^\circ L_{e_a^{(j)}}P_j^\circ
       -\frac1d I_{H_j^\circ}
 \quad\text{on }H_j^\circ.
\end{equation}
Here $P_j^\circ$ is the orthogonal projection onto $H_j^\circ$ inside
$L^2(A_j,\tau)$, and $L$ denotes left multiplication.  These formulas
also apply to the vacuum, on which only creation occurs.

Consider a fully centered word
$h_{a_r}^{(r)}\cdots h_{a_1}^{(1)}$.  The rightmost letter acts first.
Since the copy indices $1,\ldots,r$ are distinct, once a creation or
preservation occurs, every later operation must be a creation.
Thus the word is a sum of cuts of two kinds: $p$ annihilations followed
by $q=r-p$ creations, or $p$ annihilations, one preservation, and
$q=r-p-1$ creations.

For completeness, the domain restrictions in these cuts can be made
explicit.  For a set $S$ of copy indices, let $\mathcal F_{\notin S}$
denote the vacuum plus all reduced tensors whose first copy index is
outside $S$.  For a cut without preservation put
\[
 S_p=\bigl(\{p\}\text{ if }p\geq1\bigr)
       \cup\bigl(\{p+1\}\text{ if }p<r\bigr),
 \qquad E_p=\mathcal F_{\notin S_p}.
\]
Its source and target spaces are canonically embedded copies of
\[
 H_1^\circ\otimes\cdots\otimes H_p^\circ\otimes E_p,
 \qquad
 H_r^\circ\otimes\cdots\otimes H_{p+1}^\circ\otimes E_p,
\]
respectively; empty tensor products mean $\C$.  Excluding copy $p$
makes the input reduced, and excluding copy $p+1$ permits the first
creation.  For a preservation cut put $j=p+1$ and
$E=\mathcal F_{\notin\{j\}}$.  Its source and target spaces are
\[
 H_1^\circ\otimes\cdots\otimes H_p^\circ\otimes H_j^\circ\otimes E,
 \qquad
 H_r^\circ\otimes\cdots\otimes H_{j+1}^\circ\otimes H_j^\circ\otimes E.
\]
All embeddings are isometries; a cut is zero on the orthogonal
complement of its indicated source space.  In particular these restrictions
depend on the cut alone and never on the labels.

\paragraph{Distinct labels and tuple matrices.}
We now suppose that $W$ has distinct labels and sum the fully centered
words with multiset $W$.  In each cut, the annihilation map is the
tensor product of $D_1^*,\ldots,D_p^*$, with the identity on the
remaining factors.  The creation map uses the $D_j$ in decreasing
copy order.  Their norms are at most $d^{-p/2}$ and $d^{-q/2}$.
Between them is a matrix on label tuples, scalar-valued without
preservation and operator-valued on $H_j^\circ$ with preservation.
Tensor order reversals are unitary.  Together with the isometric
source and target embeddings, this gives a factorization of each cut
whose norm is bounded by
\begin{equation}\label{eq:fw-cut-factorization}
 d^{-(p+q)/2}
       \norm{\text{matrix on label tuples}}.
\end{equation}
The tail factor carries the identity throughout.

Without preservation, the tuple matrix has entry one exactly when
the input and output tuples together list $W$.  Group tuples by
their underlying subsets.  Each input subset of size $p$ is paired
only with its complementary output subset of size $q$, giving
mutually orthogonal all-ones blocks of size $q!$ by $p!$.
The matrix norm is therefore $\sqrt{p!q!}$.  Summing these cuts gives
at most
\begin{equation}\label{eq:fw-no-preservation}
 \sum_{p=0}^r\sqrt{p!(r-p)!}\,d^{-r/2}
 \leq(r+1)\sqrt{r!}\,d^{-r/2}.
\end{equation}

For a preservation cut, $p+q=r-1$.  Summing and replicating over the
orders of the input and output tuples introduces a factor
$\sqrt{p!q!}$.  The remaining matrix is indexed by subsets
$B\subset W$ of size $p$ on the input and subsets $A\subset W$ of
size $q$ on the output.  Its entry is $K_a^{(j)}$ when $A\cap B$
is empty and $\{a\}=W\setminus(A\cup B)$; other entries vanish.

To estimate this matrix, first replace $K_a^{(j)}$ by the projection
$L_{e_a^{(j)}}$ on $L^2(A_j,\tau)$ and call the resulting matrix $T$.
For a fixed row $A$, the admissible columns $B$ give distinct missing
labels $a$, so their output ranges are orthogonal.  For a fixed
column $B$, the admissible rows also give distinct $a$, so the
projection domains are orthogonal.  Thus, for an input $(\xi_B)_B$,
\[
 \norm{T\xi}^2
 =\sum_A\sum_{\substack{B:\ A\cap B=\varnothing}}
      \norm{L_{e_{a(A,B)}^{(j)}}\xi_B}^2
 \leq\sum_B\norm{\xi_B}^2.
\]
Consequently $\norm{T}\leq1$.  Compressing every input and output
leg to $H_j^\circ$ preserves this bound.  By
\eqref{eq:fw-preservation-operator}, the remaining scalar correction
is $d^{-1}$ times the zero-one matrix of admissible pairs, tensored
with $I_{H_j^\circ}$.  This matrix has $p+1$ entries in each row
and $q+1$ in each column; Cauchy--Schwarz in each row bounds its
norm by $\sqrt{(p+1)(q+1)}$.  The preservation tuple matrix therefore
has norm at most
\[
 \sqrt{p!q!}
       \left(1+\frac{\sqrt{(p+1)(q+1)}}d\right)
 \leq\sqrt{p!q!}\left(1+\frac rd\right).
\]
Using \eqref{eq:fw-cut-factorization} and summing the preservation
cuts gives
\begin{equation}\label{eq:fw-preservation-sum}
 \left(1+\frac rd\right)\sqrt{(r-1)!}\,d^{-(r-1)/2}
       \sum_{p=0}^{r-1}\binom{r-1}{p}^{-1/2}.
\end{equation}

The last sum is at most five, uniformly in $r$.  To see this, put
$m=r-1$.  For $m=0,1,2,3$ the sums are respectively
$1$, $2$, $2+1/\sqrt2$, and $2+2/\sqrt3$.  For $m\geq4$,
unimodality of the binomial coefficients gives
\[
 \sum_{p=0}^m\binom mp^{-1/2}
 \leq2+\frac2{\sqrt m}
         +(m-3)\sqrt{\frac2{m(m-1)}}
 \leq3+\sqrt2<5.
\]

\paragraph{Scalar parts and repeated labels.}
It remains to restore the scalar parts of the projections and to
treat repeated labels.  We record a bound for a single centered
word of length $l\geq1$ in distinct copy positions, with arbitrary
labels.  The same cuts now have scalar coefficients rather than
tuple sums.  Each creation or annihilation has norm at most
$d^{-1/2}$, and each preservation has norm at most one, since it is
a compression of multiplication by $e_a-d^{-1}1$.  There are $l$
preservation cuts and $l+1$ cuts without preservation.  Hence
\begin{equation}\label{eq:fw-single-word}
 \norm{h_{a_l}^{(i_l)}\cdots h_{a_1}^{(i_1)}}
 \leq l\,d^{-(l-1)/2}+(l+1)d^{-l/2}
 \leq(2l+1)d^{-(l-1)/2}.
\end{equation}
This estimate does not require the labels to be distinct.

Expand a projection word of length $r$ using $e_a=d^{-1}1+h_a$.
A term retaining $l<r$ centered letters contains the scalar factor
$d^{-(r-l)}$.  For $1\leq l<r$, \eqref{eq:fw-single-word} bounds it by
\[
 (2l+1)d^{-r+(l+1)/2}\leq(2l+1)d^{-r/2}.
\]
The term retaining no centered letters has norm $d^{-r}\leq d^{-r/2}$.
Thus the difference between $B_W$ and its fully centered sum has
norm at most
\begin{equation}\label{eq:fw-scalar-remainder}
 r! H_r d^{-r/2},
 \qquad
 H_r=1+\sum_{l=1}^{r-1}\binom rl(2l+1),
\end{equation}
because at most $r!$ ordered tuples have multiset $W$.

For distinct labels, combining
\eqref{eq:fw-no-preservation}, \eqref{eq:fw-preservation-sum}, and
\eqref{eq:fw-scalar-remainder} yields
\[
 \norm{B_W}
 \leq5\left(1+\frac rd\right)
           \sqrt{(r-1)!}\,d^{-(r-1)/2}
       +J_r d^{-r/2},
 \qquad J_r=(r+1)\sqrt{r!}+r!H_r.
\]
For example, choose an integer
\[
 d_r\geq\max\left\{5r,\frac{J_r^2}{(r-1)!}\right\}.
\]
Then the two excess terms, after division by
$\sqrt{(r-1)!}\,d^{-(r-1)/2}$, are each at most one, proving the
constant seven in \eqref{eq:fw-beta}.

For repeated labels, apply \eqref{eq:fw-single-word} to every fully
centered summand and use \eqref{eq:fw-scalar-remainder}.  Since
$d^{-r/2}\leq d^{-(r-1)/2}$, it suffices to take
\[
 C_r=r!\bigl(2r+1+H_r\bigr).
\]
These choices depend only on $r$ and complete the proof.
\end{proof}

\begin{corollary}[Sum of squared word bounds through one label]
\label{cor:through-label}
For every integer $N\geq1$ there is an integer $d_0(N)$ such that,
whenever $d\geq d_0(N)$, the constants \eqref{eq:fw-beta} satisfy
\begin{equation}\label{eq:fw-through-label}
 \max_{1\leq a\leq d}
 \sum_{\substack{W\text{ a multiset from }\{1,\ldots,d\}\\
                 |W|=r,\ a\in W}}
       \beta_W^2\leq64
 \qquad(1\leq r\leq N).
\end{equation}
Membership $a\in W$ here means that $a$ occurs at least once;
each multiset is counted once, regardless of that multiplicity.
In particular, for free copies as in \cref{lem:free-words}, the norm
bounds and \eqref{eq:fw-through-label} hold simultaneously through
length $N$ after choosing this one threshold for $d$.
\end{corollary}

\begin{proof}
Fix $r$ and a label $a$.  There are $\binom{d-1}{r-1}$ distinct-label
multisets containing $a$, and their contribution is
\[
 49(r-1)!d^{-(r-1)}\binom{d-1}{r-1}\leq49.
\]
For $r=1$ this is the entire sum.  Suppose $r\geq2$.  A repeated-label
multiset has support of cardinality $t\leq r-1$.  If it contains
$a$, its other $t-1$ support labels can be chosen in
$\binom{d-1}{t-1}$ ways.  Its positive multiplicities, with total
$r$, can be chosen in $\binom{r-1}{t-1}$ ways.  Consequently the
number of such multisets is at most
\[
 \sum_{t=1}^{r-1}\binom{d-1}{t-1}\binom{r-1}{t-1}
 \leq A_r d^{r-2},
 \qquad
 A_r=\sum_{u=0}^{r-2}\frac1{u!}\binom{r-1}{u}.
\]
Their contribution to \eqref{eq:fw-through-label} is therefore at
most $C_r^2 A_r/d$.  Choose $d_0(N)$ to be an integer at least all
the $d_r$ for $1\leq r\leq N$ and all the numbers
$C_r^2A_r/15$ for $2\leq r\leq N$.  The latter condition is empty
when $N=1$.  The sum is then at most $49+15=64$, uniformly in $a$
and simultaneously for the finitely many lengths $r\leq N$.
\end{proof}

\subsection{Free unitaries with controlled lifts}

For varying finite tracial algebras $(F_j,\tau_j)$, the same bounded-sequence construction defines the tracial ultraproduct $\prod_\omega(F_j,\tau_j)$.
A unitary $s$ is \emph{Haar} if $\tau(s^m)=0$ for every nonzero
$m\in\Z$.  Two unital subalgebras are \emph{scalar free} if every
alternating product of trace-zero elements from them has trace zero.
Subalgebras containing $B$ are \emph{free over $B$} if the corresponding
assertion holds with the trace replaced by the trace-preserving
conditional expectation $E_B$.

\begin{lemma}[Free Haar unitaries with matrix-algebra lifts]
\label{lem:free-lifts}
Let $M\subseteq\B(H)$ be a separably acting $\II_1$ factor, let
$n\geq1$, and put $P=M_n(M)$.  For every free ultrafilter $\omega$
there is a Haar unitary $s\in P^\omega$ such that $W^*(s)$ and the
constant copy of $P$ are scalar free.  Moreover, $s$ has unitary
representatives
\[
 s=(s_j)_\omega,\qquad
 s_j\in\U(M_n(R_j)),
\]
where each $R_j\subseteq M$ is a unital full matrix algebra.
Consequently, for every $h=h^*\in\B(H)$, with
$\delta=\norm{\ad_h|_M}$ and $h_n=I_n\otimes h$, these representatives
satisfy
\begin{equation}\label{eq:free-lift-commutator}
 \norm{[h_n,s_j]}\leq2\delta\qquad(j\geq1).
\end{equation}
\end{lemma}

\begin{proof}
We construct the free unitary in two stages, then control its representatives.

\paragraph{Ultraproduct facts and the independence input.}
We first record the two facts about tracial ultraproducts used in the
construction.  A tracial ultraproduct of finite factors is a factor.
Indeed, in any finite factor $F$ the scalar $\tau(x)1$ belongs to the
$2$-norm closed convex hull of $\{uxu^*:u\in\U(F)\}$.  To see this,
take the unique element of least $2$-norm in that closed convex hull.
It is fixed by every unitary conjugation.  It still belongs to $F$,
because the convex hull is uniformly bounded in operator norm and
the same is true of its $2$-norm closure.  It is therefore scalar,
and its trace identifies it as $\tau(x)1$.  In particular,
\[
 \norm{x-\tau(x)1}_2
 \leq\sup_{u\in\U(F)}\norm{[x,u]}_2.
\]
If $(x_j)_\omega$ is central in an ultraproduct, the supremum on the
right tends to zero along $\omega$: otherwise choosing an almost
maximizing unitary in each offending coordinate would give a unitary
sequence not commuting with $(x_j)_\omega$.  Thus every central
element is scalar.  The ultraproduct of matrix factors whose sizes
tend to infinity is diffuse, since their projections have traces
approaching any prescribed number in $[0,1]$.  An ultrapower of a
$\II_1$ factor is likewise diffuse, since it contains the constant
factor.  These ultraproducts are therefore $\II_1$ factors.

Second, every unitary in such an ultraproduct lifts to coordinate
unitaries.  If $(x_j)_\omega$ is unitary, then
$\lim_\omega\norm{x_j^*x_j-1}_2=0$.  Extend the partial isometry in
the polar decomposition of $x_j$ to a unitary $u_j$.  This extension
exists because the two complementary support projections have the
same trace in the finite factor.  Functional calculus gives
\[
 \norm{x_j-u_j}_2
 =\norm{|x_j|-1}_2
 \leq\norm{x_j^*x_j-1}_2,
\]
so $(u_j)_\omega=(x_j)_\omega$.

We use the following form of
\citet[Theorem~0.1(b)]{Popa2014}.  Let
$\mathcal F=\prod_\omega F_j$, where the $F_j$ are finite factors
with $\dim F_j\longrightarrow\infty$, allowing infinite dimensions.
For a separable amenable von Neumann subalgebra $B_1\subseteq\mathcal F$,
put $Q=B_1'\cap\mathcal F$ and $D=Q'\cap\mathcal F$.
If $\mathcal X\subseteq\ker E_D$ is a separable linear subspace,
there is a diffuse abelian von Neumann algebra $A\subseteq Q$ such that
\begin{equation}\label{eq:popa-exact-endpoints}
 E_D(x_0a_1x_1\cdots a_rx_r)=0
\end{equation}
whenever $r\geq1$, $x_0\in\mathcal X\cup\{1\}$,
$x_1,\ldots,x_r\in\mathcal X$, and
$a_1,\ldots,a_r\in A\cap\ker\tau$.
Here separability is for the $2$-norm, as specified in
\cite[Section~1.1]{Popa2014}.  We use the displayed form with its
final $x_r$ and establish the additional endpoints needed below.

In our applications $D$ is a full matrix algebra,
$A\subseteq D'\cap\mathcal F$, and $\mathcal X$ is selfadjoint.
For $a\in A$, bimodularity implies $E_D(a)\in Z(D)$, whence
$E_D(a)=\tau(a)1$.  Taking adjoints in
\eqref{eq:popa-exact-endpoints} also treats alternating words that
start with an $\mathcal X$-letter and end with an $A$-letter.
For a word with both endpoints in $A$, write
\[
 W=a_0x_1a_1\cdots x_ra_r,
 \qquad r\geq1,
\]
with all displayed letters centered in their respective senses.
For $b\in D$, traciality and $[A,D]=0$ give
\[
 \tau(bW)=\tau\bigl(b\,a_ra_0x_1a_1\cdots x_r\bigr).
\]
Split $a_ra_0$ into its trace and a trace-zero element of $A$.
The latter term vanishes by \eqref{eq:popa-exact-endpoints} with
$x_0=1$; the scalar term vanishes by the same formula when $r\geq2$,
and by $E_D(x_1)=0$ when $r=1$.  Thus $E_D(W)=0$.
A single centered $A$-letter has already been covered.  Hence all
alternating endpoint patterns used below are valid.

\paragraph{A finite-matrix model.}
We first construct a matrix-algebra model that is free from
$B=M_n(\C)\otimes1$.  In
\[
 \mathcal F_0=\prod_\omega\bigl(M_n(\C)\otimes M_j(\C)\bigr)
\]
apply Popa's theorem with $B_1=\C1$, so that $Q=\mathcal F_0$ and
$D=\C1$, and with $\mathcal X=B\cap\ker\tau$.
The resulting diffuse abelian algebra is scalar free from $B$ by
the endpoint argument.  It contains a Haar unitary $\widetilde s_0$:
successively bisect projections in this diffuse abelian algebra to
obtain a selfadjoint element uniformly distributed on $[0,1]$, and
exponentiate it by $t\mapsto e^{2\pi it}$.
Choose unitary lifts
$\widetilde s_{0,j}\in M_n(\C)\otimes M_j(\C)$.

For each $j$, choose a unital matrix subalgebra
$R_j^0\cong M_j(\C)$ of $M$.  Such a subalgebra is obtained by
subdividing $1$ into $j$ projections of trace $1/j$ and choosing
matrix units using equivalence of equal-trace projections in a
finite factor.  The coordinate embeddings
$M_n(\C)\otimes M_j(\C)\hookrightarrow M_n(M)$ preserve the
normalized traces and fix $B$.  They induce a trace-preserving
embedding of $\mathcal F_0$ into
\[
 \mathcal P=P^\omega=M_n(M^\omega).
\]
Thus the image $s_0$ of $\widetilde s_0$ is Haar and scalar free
from $B$, and it has representatives
$s_{0,j}\in\U(M_n(R_j^0))$.

\paragraph{Freeness from the constant algebra.}
Next we move this model into free position relative to the whole
constant algebra $P$.  The equality
$P^\omega=M_n(M^\omega)$ follows entry by entry from
$\norm{(x_{ab})}_2^2=n^{-1}\sum_{a,b}\norm{x_{ab}}_2^2$.
Since $M^\omega$ is a factor,
\[
 B'\cap\mathcal P=I_n\otimes M^\omega,
 \qquad (B'\cap\mathcal P)'\cap\mathcal P=B.
\]
Put $T=W^*(P,s_0)\subseteq\mathcal P$ and
$\mathcal X=T\cap\ker E_B$.  This is a $2$-norm separable
selfadjoint linear subspace: the separably acting algebra $P$ is
countably generated as a von Neumann algebra, and adjoining one
unitary preserves this property.  Equivalently, bounded polynomial
approximations in a countable generating set give $2$-norm density.
Apply Popa's theorem with $B_1=B$ and this $\mathcal X$.
Choose a Haar unitary $w$ in the resulting abelian algebra
$A\subseteq I_n\otimes M^\omega$.

Let $P_1=W^*(B,s_0)$ and $P_2=wP_1w^*$.
Both contain $B$, since $w$ commutes with $B$.  Trace testing against
$B$ shows that
\[
 E_B(wyw^*)=E_B(y)\qquad(y\in\mathcal P).
\]
We claim that $P_2$ and $P$ are free over $B$.  In an alternating
product of $B$-centered elements from these algebras, write each
$P_2$-letter as $wyw^*$ with $y\in P_1\cap\ker E_B$.
The resulting word alternates letters $w,w^*$ from $A\cap\ker\tau$
with letters in $\mathcal X$.  It has zero $E_B$ by
\eqref{eq:popa-exact-endpoints} and the endpoint extensions proved
above.  A product consisting of one letter has zero expectation
by its assumed centering.  This proves the claim.

Set $s=ws_0w^*$ and $C_s=W^*(s)$.
It is Haar, and $C_s$ is scalar free from $B$, because conjugation
by $w$ fixes $B$ and preserves the trace.  To deduce scalar
freeness from $P$, first observe that
\begin{equation}\label{eq:scalar-free-block}
 E_B(c_1b_1c_2\cdots b_{m-1}c_m)=0
\end{equation}
for $m\geq1$, $c_i\in C_s\cap\ker\tau$, and
$b_i\in B\cap\ker\tau$.  Indeed, test the trace against an
arbitrary $b\in B$ and split $b$ into its scalar and centered
parts; scalar freeness makes both resulting traces zero.
By bimodularity, arbitrary $B$-letters may also be attached at the
two endpoints in \eqref{eq:scalar-free-block}.

Single letters have trace zero by centering.  Now take any mixed
alternating scalar-centered word from $C_s$ and $P$,
and split each $P$-letter $x$ as
\[
 x=(x-E_B(x))+E_B(x).
\]
The first summand is $B$-centered, and the second lies in
$B\cap\ker\tau$.  In every expanded term, group maximal blocks
of $C_s$-letters and the selected $B$-letters.  Each such block
lies in $P_2$ and is $B$-centered by
\eqref{eq:scalar-free-block}, including any $B$-letters at the
endpoints.  These blocks alternate with $B$-centered $P$-letters,
so freeness over $B$ makes their product have zero expectation.
If there is no remaining $B$-centered $P$-letter, the term is one
block covered by \eqref{eq:scalar-free-block}.  Taking traces
proves that $C_s$ and $P$ are scalar free.

\paragraph{Controlled representatives.}
Finally write $w=I_n\otimes z$ with $z\in\U(M^\omega)$, and
choose unitary lifts $z_j\in\U(M)$.  Then
\[
 R_j=z_jR_j^0z_j^*,\qquad
 s_j=(I_n\otimes z_j)s_{0,j}(I_n\otimes z_j^*)
\]
have all the asserted properties.  For the commutator estimate,
average $h$ over the compact unitary group of $R_j$:
\[
 h^{(j)}=\int_{\U(R_j)}u h u^*\,du.
\]
This operator commutes with $R_j$, and
$\norm{h-h^{(j)}}\leq\delta$.  Therefore
$I_n\otimes h^{(j)}$ commutes with $s_j$, and
\[
 \norm{[h_n,s_j]}
 =\norm{[I_n\otimes(h-h^{(j)}),s_j]}
 \leq2\norm{h-h^{(j)}}\leq2\delta.
\]
\end{proof}

\begin{lemma}[Unitary twists of a free Haar unitary]
\label{lem:haar-twist}
Let $(F,\tau)$ be a finite tracial von Neumann algebra, let
$P\subseteq F$ be a unital von Neumann subalgebra, and let $s\in F$
be Haar and scalar free from $P$.  For every $U\in\U(P)$, both
$sU$ and $Us$ are Haar and scalar free from $P$.
\end{lemma}

\begin{proof}
Put $v=sU$.  It suffices first to show
\[
 \tau(v^k)=0\quad(k\neq0),\qquad
 \tau(v^{k_1}x_1\cdots v^{k_r}x_r)=0
\]
for $r\geq1$, nonzero integers $k_i$, and
$x_i\in P\cap\ker\tau$.  These cyclic mixed words suffice:
if an alternating word starts and ends in the same algebra,
cyclically bring those two letters together, split their product
into its scalar and centered parts, and induct on word length.
Regard an expanded trace word as a cyclic
sequence of letters $s$ and $s^*$ separated by elements of $P$.
The expansion of $(sU)^k$ uses $(sU)^k$ itself for $k>0$ and
$(U^*s^*)^{-k}$ for $k<0$.  At a transition from an $s$-letter to
an $s^*$-letter the separator is $Ux_iU^*$; at a transition from
$s^*$ to $s$ it is $x_i$.  Thus every separator between opposite
signs is centered, including the cyclic junction.

Split every other separator into its scalar and centered parts.
In each resulting term, removal of a scalar separator merges only
powers of $s$ with the same sign.  Consequently every surviving
power has nonzero exponent, and every surviving separator is a
centered element of $P$.  If any separator remains, freeness of
$s$ from $P$ gives trace zero.  If no separator remains, all signs
were the same and the term is a scalar multiple of a nonzero power
of $s$, again with trace zero.  The argument for $\tau(v^k)$ is
identical, with all signs the same from the outset.
This proves the displayed identities.

Centered Laurent polynomials in $v$ are $2$-norm dense, with
bounded approximants, in the centered part of $W^*(v)$: use the
Fej\'er sums of the corresponding bounded functions on the circle.
The displayed identities therefore imply scalar freeness of
$W^*(v)$ and $P$.  Here bounded $2$-norm approximation can be
passed through each trace word by telescoping and the inequality
$\abs{\tau(a y b)}\leq\norm{a}\norm{y}_2\norm{b}$.
Finally apply the proved assertion to $s^*$ and $U^*$ and take
adjoints to obtain the assertion for $Us$.
\end{proof}

\subsection{Randomized coefficients}

We now prove \Cref{thm:amplification}. Fix \(M\subseteq\B(H)\) and
selfadjoint \(h\), and set \(\delta=\norm{\ad_h|_M}\).
We may assume \(\delta>0\). It suffices first to bound the
commutator with an arbitrary unitary at an arbitrary matrix level.
Fix \(n\ge2\), \(U\in\U(P)\), where \(P=M_n(M)\), and a unit vector
\(X\in H^n\). Write
\[
 h_n=I_n\otimes h,\qquad Y=UX,\qquad
 B=M_n(\C)\otimes1.
\]
The following order of choices will be used throughout this proof:
\[
 R_*=100\pi,\qquad
 N\ge100nR_*^3,\qquad
 k\ge N,\qquad d=nk.
 \tag{\(\ast\)}\label{eq:amplification-choices}
\]
Here \(N\) is an integer, and then \(k\) is chosen so large that
\Cref{cor:through-label} holds at every length \(1\le r\le N\).
Choose a unital matrix algebra \(R\cong M_k(\C)\) in \(M\).
Such an algebra is obtained by subdividing the identity into \(k\)
equivalent projections of trace \(1/k\) and choosing matrix units.
Let \(e_1,\ldots,e_d\) be the diagonal minimal projections of
\(D=M_n(R)\).

Take the free unitary \(s=(s_j)_\omega\) of \Cref{lem:free-lifts}.
Each \(s_j\) lies in \(M_n(R_{0,j})\) for a unital matrix algebra
\(R_{0,j}\subset M\), and therefore
\begin{equation}\label{eq:cheap-lift-commutator}
 \norm{[h_n,s_j]}\le2\delta.
\end{equation}
Put \(v_j=s_jU\) and \(v=sU\). By \Cref{lem:haar-twist}, \(v\)
is a Haar unitary free from the constant algebra \(P\).

Fix an integer \(q>2N\). Independently and uniformly choose
\(q\)-th roots of unity \(\theta_1,\ldots,\theta_d\), and set
\[
 t=\sum_{a=1}^d\theta_ae_a,\qquad L_j=tv_j.
\]
For \(1\le r\le N\), consider the unitary expressions
\[
 Q^+_{r,j}=(tv_j)^{r-1}t,\qquad
 Q^-_{r,j}=(v_j^*t^*)^r.
\]
In either expression, group its phase monomials by the multiset
\(W\) of the \(r\) labels that occur. Thus
\[
 Q^\pm_{r,j}=\sum_{|W|=r}\chi^\pm_W(\theta)K^\pm_{W,j},
 \qquad
 \chi^\pm_W(\theta)=\prod_{a\in W}\theta_a^{\pm1}.
\]
Multiplicities are included in the product. Within either sign, these
characters are distinct and orthogonal, since \(q>2N\).
For the next estimates we suppress the superscript and \(r\).
Orthogonality and unitarity give
\begin{equation}\label{eq:coefficient-column}
 \sum_WK_{W,j}^*K_{W,j}=I.
\end{equation}
Moreover,
\begin{equation}\label{eq:coefficient-support}
 K_{W,j}=K_{W,j}p_W,\qquad
 p_W=\sum_{a\in W\text{ distinct}}e_a,
\end{equation}
because the rightmost factor carrying a phase is a diagonal
projection with a label in \(W\).

Their classes \(K_W=(K_{W,j})_\omega\) have norm at most
\(\beta_W\) from \Cref{lem:free-words}.
Indeed, moving all powers of \(v\) to the left rewrites each
coefficient, up to a left unitary, as the corresponding sum of
projection words from distinct algebras \(v^{-l}Dv^l\).
These algebras are scalar free: in a product of centered letters
from alternating conjugates, consecutive unequal exponents leave
nonzero powers of the Haar unitary between centered \(P\)-letters;
scalar freeness of \(v\) and \(P\) makes its trace zero.
Consequently, \Cref{cor:through-label} gives
\begin{equation}\label{eq:beta-through-label}
 \max_a\sum_{\substack{|W|=r\\a\in W}}\beta_W^2\le64
 \qquad(1\le r\le N).
\end{equation}

The coordinate coefficients need not themselves obey their limiting
norm bounds. We enforce those bounds by spectral truncation. For
\(b>0\), put
\[
 f_b(u)=\min\{1,b/u\}\quad(u>0),\qquad f_b(0)=1,
\]
and define
\[
 K'_{W,j}=K_{W,j}f_{\beta_W}(|K_{W,j}|),\qquad
 E_{W,j}=K_{W,j}-K'_{W,j}.
\]
Then \(\norm{K'_{W,j}}\le\beta_W\), its right support is still
contained in \(p_W\), and
\begin{equation}\label{eq:clipped-columns}
 \sum_W(K'_{W,j})^*K'_{W,j}\le I,\qquad
 \sum_WE_{W,j}^*E_{W,j}\le I.
\end{equation}
These follow term by term from \Cref{eq:coefficient-column}.
Continuous functional calculus in the quotient gives
\[
 \lim_{j\to\omega}\norm{E_{W,j}}_2=0,
\]
since the positive function \((u-\beta_W)_+\) vanishes on
the spectrum of \(|K_W|\).

The purpose of these truncated coefficients is to replace a matrix
commutator by a column commutator with a small support. Average \(h\)
over \(\U(R)\) with Haar probability measure, obtaining
\[
 h^0=\int_{\U(R)}uhu^*\,du\in R',\qquad
 \norm{h-h^0}\le\delta,\qquad
 \norm{\ad_{h^0}|_M}\le\delta.
\]
Write \(h_n^0=I_n\otimes h^0\).
The last inequality follows by averaging
\([uhu^*,a]=u[h,u^*au]u^*\).
For each \(W\), there is a rectangular partial isometry
\(V\in M_{n,1}(R)\) with \(VV^*=p_W\): the support of \(p_W\)
has at most \(r\le k\) minimal diagonal projections, so its rank
does not exceed the dimension \(k\) of the source column.
Its entries commute with \(h^0\), and hence
\[
 [h_n^0,K'_{W,j}]Z
 =\bigl(h_n^0K'_{W,j}V-K'_{W,j}Vh^0\bigr)V^*Z.
\]
\Cref{lem:column-commutator} bounds this by
\(8\delta\beta_W\norm{p_WZ}\). Orthogonality of the diagonal
projections and \Cref{eq:beta-through-label} give the weighted support
estimate
\[
 \sum_W\beta_W^2\norm{p_WZ}^2
 =\sum_{a=1}^d\left(\sum_{\substack{|W|=r\\a\in W}}\beta_W^2\right)
       \norm{e_aZ}^2
 \le64\norm{Z}^2.
\]
Squaring the column bounds, summing over \(W\), and taking a square
root therefore yields
\[
 \left(\sum_W\norm{[h_n^0,K'_{W,j}]Z}^2\right)^{1/2}
 \le64\delta\norm{Z}.
\]
Replacing \(h^0\) by \(h\) costs at most \(2\delta\norm{Z}\),
by the first column contraction in \Cref{eq:clipped-columns}.
Thus
\begin{equation}\label{eq:clipped-drift}
 \left(\sum_W\norm{[h_n,K'_{W,j}]Z}^2\right)^{1/2}
 \le66\delta\norm{Z}.
\end{equation}

For \(Z=X\), restoring the coefficients \(E_{W,j}\) costs a quantity
tending to zero along \(\omega\), by \Cref{lem:trace-transfer}.
There are only finitely many \(W\).
For the moving vector \(Z=v_jX=s_jY\), write
\[
 [h_n,E_{W,j}]s_jY
 =h_nE_{W,j}s_jY-E_{W,j}s_jh_nY
      -E_{W,j}[h_n,s_j]Y.
\]
The first two terms tend to zero in column norm: \(E_{W,j}s_j\)
is bounded and 2-small, and \(Y,h_nY\) are fixed.
The last term has column norm at most \(2\delta\), by
\Cref{eq:clipped-columns,eq:cheap-lift-commutator}.
Phase orthogonality now converts these column bounds back to
root-mean-square bounds for the unitary words:
\begin{align}
 \lim_{j\to\omega}
 \left(\E_\theta\norm{[h_n,Q^+_{r,j}]v_jX}^2\right)^{1/2}
 &\le68\delta,\label{eq:positive-drift}\\
 \lim_{j\to\omega}
 \left(\E_\theta\norm{[h_n,Q^-_{r,j}]X}^2\right)^{1/2}
 &\le66\delta.\label{eq:negative-drift}
\end{align}
Both bounds hold for every \(r\le N\).

\subsection{Recovering a scalar operator}

The word estimates control commutator errors independently of the matrix
size. We now combine the words into rows and choose a contraction of
$M$ so that their averaged product converges weakly to a positive scalar
multiple of the identity. This scalar will detect the matrix commutator.

Independently choose a uniform \(q\)-th root of unity \(z\) and
a uniform standard basis column \(\alpha\in\C^n\).
Write \(p=\alpha\alpha^*\in B\), and define
\[
 f_j=N^{-1/2}\sum_{r=1}^N(zL_j)^r,\qquad
 F_j=\sqrt n\,\alpha^*f_j,\qquad
 G_j=\sqrt n\,\alpha^*f_j^*.
\]
Here \(F_j,G_j:H^n\to H\) are rows over \(M\).
In the remainder of the proof \(\E\) averages over all three finite
random choices \(\theta,z,\alpha\). Averaging first over \(\alpha\)
and then using orthogonality of \(z,\ldots,z^N\) gives
\begin{equation}\label{eq:row-rms}
 \E\norm{F_jX}^2=\E\norm{G_jX}^2=1.
\end{equation}

The same orthogonality removes cross terms between different powers
in the following errors:
\begin{align}
 \lim_{j\to\omega}
 \left(\E\norm{hF_jX-F_jU^*h_nY}^2\right)^{1/2}
 &\le70\delta,\label{eq:forward-row-drift}\\
 \lim_{j\to\omega}
 \left(\E\norm{hG_jX-G_jh_nX}^2\right)^{1/2}
 &\le70\delta.\label{eq:backward-row-drift}
\end{align}
To check the first estimate, use \(L_j^r=Q^+_{r,j}s_jU\).
The corresponding error before the scalar phase is
\[
 [h_n,Q^+_{r,j}]s_jY+Q^+_{r,j}[h_n,s_j]Y,
\]
which is bounded by \Cref{eq:positive-drift} and
\Cref{eq:cheap-lift-commutator}. For the second estimate,
the adjoint power is \(Q^-_{r,j}\), so
\Cref{eq:negative-drift} applies. Averaging the normalized sum of
the \(N\) squared bounds produces no factor depending on \(N\).

We next choose contractions \(x_j\in M\) so that the average of
\(G_j^*x_jF_j\) has a nonzero scalar weak limit. Put
\[
 A_j=pf_j^2p,\qquad
 \widetilde x_j=A_j^*(A_jA_j^*+.01p)^{-1/2}
              =\alpha x_j\alpha^*.
\]
The inverse square root is taken in the unital corner \(pPp\).
Functional calculus gives \(\norm{x_j}\le1\). Let
\[
 C_j=G_j^*x_jF_j=nf_j\widetilde x_jf_j.
\]
Since \(\norm{f_j}\le\sqrt N\), these operators have norm at
most \(nN\), uniformly in \(j\) and the random choices.

Use the same letters without the subscript for their ultrapower
classes. The unitary \(u=ztv\) is Haar and free from \(P\), by
\Cref{lem:haar-twist}. Since \(C\in W^*(u,p)\), freeness gives
\begin{equation}\label{eq:scalar-expectation}
 E_P(C)=\lambda p+\mu(1-p).
\end{equation}
For completeness, expand a polynomial in \(u,u^*,p\) into alternating
centered letters from \(W^*(u)\) and \(W^*(p)\).
Every reduced term containing a centered \(u\)-letter has zero
expectation onto \(P\), as is seen by testing its trace against
elements of \(P\). The remaining terms lie in \(W^*(p)\).
Polynomial approximation, including the continuous corner functional
calculus defining \(\widetilde x\), proves
\Cref{eq:scalar-expectation}. The numbers \(\lambda,\mu\)
depend only on \(n,N\): the free joint law of \(u,p\) is the same
for all choices, since \(\tau(p)=1/n\).
As \(\E_\alpha p=I/n\), it follows that
\begin{equation}\label{eq:scalar-average}
 \E E_P(C)=c_*I,\qquad c_*=\tau(C).
\end{equation}

We give the lower bound on \(c_*\), since its independence from \(n\)
is essential. Traciality and the support of the corner functional
calculus show that
\[
 c_*=n\tau\bigl(AA^*(AA^*+.01p)^{-1/2}\bigr)
 \ge n\norm{pf^2p}_1-.1.
\]
The inequality follows from
\(t^2(t^2+.01)^{-1/2}\ge t-.1\) for \(t\ge0\), and
\(n\tau(p)=1\).

Let \(e\) be the spectral projection of \(u\) for the arc
\(\{e^{it}:|t|\le R_*/N\}\), and put
\(\eta=\tau(e)=R_*/(\pi N)\).
Haar measure gives \(\tau(|f|^2)=1\), and
\(\norm{f^2}\le N\).
For \(0<|t|\le\pi\), the geometric-sum formula gives
\[
 \left|\sum_{r=1}^Ne^{irt}\right|\le\frac{\pi}{|t|}.
\]
Integrating its square outside the arc yields
\begin{equation}\label{eq:haar-tail}
 \tau(|f|^2(1-e))
 \le \frac1{2\pi}\int_{|t|>R_*/N}
               \frac{\pi^2}{Nt^2}\,dt
 \le\frac{\pi}{R_*}.
\end{equation}
Since \(f\) commutes with \(e\), tracial Cauchy--Schwarz and
freeness of \(p\) from \(f,e\) give
\begin{align*}
 n\norm{p(f^2-ef^2e)p}_1
 &=n\norm{pf(1-e)fp}_1\\
 &\le n\norm{pf(1-e)}_2\norm{(1-e)fp}_2\\
 &=\tau(|f|^2(1-e))
 \le\frac{\pi}{R_*}.
\end{align*}

Put \(s_e=nepe\), a positive operator in the corner \(eP^\omega e\).
The polar decomposition of \(\sqrt n\,ep\) identifies the nonzero
parts of the following two operators, and hence gives
\begin{equation}\label{eq:compression-polar}
 n\norm{pef^2ep}_1
 =\norm{s_e^{1/2}f^2s_e^{1/2}}_1.
\end{equation}
Explicitly, if \(\sqrt n\,ep=s_e^{1/2}w\), the left-hand operator
inside the norm is
\(w^*s_e^{1/2}f^2s_e^{1/2}w\), and the middle operator is supported
on \(ww^*=\supp(s_e)\) on both sides.

The elementary free moment calculation gives
\[
 \norm{s_e^{1/2}}_2=\sqrt\eta,\qquad
 \norm{s_e-e}_2^2=(n-1)\eta^2.
\]
For the second identity, put \(a=np-I\). Then
\(\tau(a)=0\), \(\tau(a^2)=n-1\), and freeness gives
\(\tau(aeae)=\eta^2\tau(a^2)\); expansion of
\(\tau((nepe-e)^2)\) gives the same quantity.
Within the corner with identity \(e\),
\(|\sqrt t-1|\le|t-1|\) for \(t\ge0\), so
\(\norm{s_e^{1/2}-e}_2\le\sqrt{n-1}\eta\).
Consequently,
\begin{align*}
 \norm{s_e^{1/2}f^2s_e^{1/2}-ef^2e}_1
 &\le\norm{s_e^{1/2}-e}_2\norm{f^2}
       \bigl(\norm{s_e^{1/2}}_2+\norm{e}_2\bigr)\\
 &\le2N\sqrt{n-1}\eta^{3/2}.
\end{align*}
Since \(\norm{ef^2e}_1=\tau(e|f|^2)\), we obtain
\begin{equation}\label{eq:positive-scalar}
 c_*\ge .9-\frac{2\pi}{R_*}
       -2N\sqrt{n-1}\left(\frac{R_*}{\pi N}\right)^{3/2}
 \ge\frac12.
\end{equation}
Indeed, by \Cref{eq:amplification-choices}, the last term is at most
\(.2/\pi^{3/2}\), and \(2\pi/R_*=.02\).

For each fixed random choice, \Cref{lem:trace-transfer} identifies
the weak ultralimit of \(C_j\) with \(E_P(C)\).
All random spaces are finite, and their sizes are fixed before
taking the ultralimit. Thus \Cref{eq:scalar-average} gives
\begin{equation}\label{eq:weak-scalar-limit}
 \underset{j\to\omega}{\operatorname{WOT}\!-\!\lim}\,\E C_j=c_*I.
\end{equation}

\subsection{Extracting the matrix commutator}

The randomized rows now have two properties: their commutator drift
is bounded independently of \(n\), and their product with a
contraction \(x_j\in M\) has a positive scalar weak limit.
Apply the defining scalar commutator bound to that contraction:
\[
 \left|\E\ip{[h,x_j]F_jX}{G_jX}\right|
 \le\delta\,\E\bigl(\norm{F_jX}\norm{G_jX}\bigr)\le\delta,
\]
using \Cref{eq:row-rms}.
As \(h=h^*\), the expression inside the average equals
\[
 \ip{x_jF_jX}{hG_jX}-\ip{x_jhF_jX}{G_jX}.
\]
Substitute the main terms from
\Cref{eq:forward-row-drift,eq:backward-row-drift}.
By Cauchy--Schwarz and \Cref{eq:row-rms}, the sum of the two
average errors has ultralimit at most \(140\delta\).
The remaining expression is
\[
 \ip{C_jX}{h_nX}-\ip{C_jU^*h_nY}{X}.
\]
\Cref{eq:weak-scalar-limit} therefore implies
\[
 c_*\left|\ip{h_nX}{X}-\ip{h_nUX}{UX}\right|
 \le141\delta.
\]
Here the quadratic forms are real because \(h_n\) is selfadjoint.
By \Cref{eq:positive-scalar}, their difference has absolute value
at most \(282\delta\). The choices of randomization were allowed
to depend on \(U,X,n\), but this bound does not.
Taking the supremum over unit vectors gives
\[
 \norm{h_n-U^*h_nU}\le282\delta,\qquad
 \norm{[h_n,U]}\le282\delta.
\]
Every selfadjoint contraction \(a\) in a unital \(C^*\)-algebra is
the average of the unitaries
\(a\pm i(1-a^2)^{1/2}\).
Writing a general contraction as its real part plus \(i\) times
its imaginary part therefore bounds its commutator by
\(564\delta\).
The level \(n=1\) is already covered by the definition of \(\delta\).
Taking the supremum over \(n\) proves
\Cref{thm:amplification}, for example with
\(C_{\mathrm{amp}}=600\).
